\unitopen{10}
\begin{candos}
\item opisać swoje mieszkanie i jego pokoje;
\item powiedzieć, gdzie coś stoi lub leży (\de{Wo?}), i dokąd idziesz (\de{Wohin?});
\item użyć przyimków lokalnych z odpowiednim przypadkiem;
\item zapytać o mieszkanie do wynajęcia: czynsz, koszty dodatkowe, kaucja, termin oglądania.
\end{candos}

\section{Słownictwo}
\vocabtable{10}

\section{Przyimki lokalne: \de{Wo?} i \de{Wohin?}}
Dziewięć przyimków (\de{in, auf, an, unter, über, neben, vor, hinter, zwischen}) łączy się z \emph{dwoma} przypadkami:
\textbf{celownik} (\de{Dativ}) odpowiada na \de{Wo?} (miejsce, brak ruchu do celu), \textbf{biernik} (\de{Akkusativ}) na
\de{Wohin?} (ruch do celu). Biernik już znasz z jednostki~5. Celownik ma inne rodzajniki:

\begin{center}\small
\begin{tabular}{@{}lllll@{}}
\toprule
 & \textsf{\textbf{męski}} & \textsf{\textbf{żeński}} & \textsf{\textbf{nijaki}} & \textsf{\textbf{mnoga}}\\
\midrule
\de{Wo?} (Dativ) & \artm{dem} / \de{einem} & \artf{der} / \de{einer} & \artn{dem} / \de{einem} & \artf{den} + \de{-n}\\
\de{Wohin?} (Akkusativ) & \artm{den} / \de{einen} & \artf{die} / \de{eine} & \artn{das} / \de{ein} & \artf{die}\\
\bottomrule
\end{tabular}
\end{center}
Zlewanie: \de{in dem = im}, \de{an dem = am}, \de{in das = ins}, \de{an das = ans}, \de{auf das = aufs}.

\de{Der Laptop liegt \textbf{auf dem} Tisch.} (Wo?) $\rightarrow$ \de{Ich lege den Laptop \textbf{auf den} Tisch.} (Wohin?)
\de{Ich bin \textbf{im} Büro.} $\rightarrow$ \de{Ich gehe \textbf{ins} Büro.} \de{Wir wohnen \textbf{in der} Stadt.} $\rightarrow$
\de{Wir fahren \textbf{in die} Stadt.}

\begin{zasada}[Dwie pary czasowników]
\textbf{Stan:} \de{stehen}, \de{liegen}, \de{sitzen}, \de{hängen} $+$ celownik (\de{Wo?}). \textbf{Ruch}: \de{stellen}, \de{legen},
\de{setzen}, \de{hängen} $+$ biernik (\de{Wohin?}). \de{Die Flasche steht auf dem Tisch. Ich stelle die Flasche auf den Tisch.}
\end{zasada}

\begin{falle}
Po polsku „na stole” i „na stół” różnią się końcówką, a po niemiecku zmienia się \emph{rodzajnik} i
\emph{sposób myślenia}: pytasz sam siebie \de{Wo?} czy \de{Wohin?}. Typowy błąd: \de{Ich bin in das Büro} (gdy już jesteś
w biurze). Poprawnie: \de{Ich bin im Büro.} Drugi: \de{Ich gehe zu Hause} (dosłownie „idę w domu”). Poprawnie:
\de{Ich gehe nach Hause} (dokąd?) i \de{Ich bin zu Hause} (gdzie?).
\end{falle}

\section{Szukanie mieszkania}
Ogłoszenie czyta się w skrótach: \de{3-Zi-Whg} (trzypokojowe mieszkanie), \de{EG} (parter), \de{OG} (piętro), \de{inkl.\ NK}
(z kosztami dodatkowymi), \de{ab sofort} (od zaraz). Pytania: \de{Wie hoch ist die Miete?} \de{Was sind die Nebenkosten?}
\de{Gibt es eine Kaution?} \de{Wann kann ich die Wohnung besichtigen?}

\begin{szwajcaria}
W Szwajcarii mieszkanie „3,5 Zimmer” oznacza zwykle trzy pokoje i dodatkowy pół-pokój (np.\ kuchnia/przedpokój)
\verify{(konwencja liczenia bywa różna, pytaj o metraż w m\textsuperscript{2})}. Rynek najmu w Zurychu jest bardzo napięty, a kaucja wynosi
zwykle do kilku miesięcznych czynszów \verify{(górna granica jest określona prawem, sprawdź w OR i u wynajmującego)}.
Szwajcaria używa \de{Mietzins} jako słowa oficjalnego, a dokumenty: \de{Mietvertrag}, \de{Betreibungsauszug}
(zaświadczenie o postępowaniach egzekucyjnych, którego chcą wynajmujący). Koszty dodatkowe to \de{Nebenkosten}.
\end{szwajcaria}

\begin{praca}
Home office i hybryda: \de{Ich arbeite heute von zu Hause.} \de{Ich bin im Homeoffice.} \de{Ich komme morgen ins Büro.}
\de{Der Server steht im Rechenzentrum.} Zwróć uwagę na \de{von zu Hause} (z domu) oraz \de{im Büro} / \de{ins Büro}.
\end{praca}

\section{Dialogi}
\textbf{Dialog 1 --- telefon w sprawie mieszkania}
\begin{dialog}
\turn{A}{Guten Tag, ich rufe wegen der Wohnung in der Bahnhofstraße an.}{Dzień dobry, dzwonię w sprawie mieszkania przy Bahnhofstraße.}
\turn{B}{Ja, die Wohnung ist noch frei. Sie hat drei Zimmer, eine Küche und ein Bad.}{Tak, mieszkanie jest jeszcze wolne. Ma trzy pokoje, kuchnię i łazienkę.}
\turn{A}{Wie hoch ist die Miete?}{Jak wysoki jest czynsz?}
\turn{B}{Die Miete beträgt 1.900 Franken plus 200 Franken Nebenkosten.}{Czynsz wynosi 1900 franków plus 200 franków kosztów dodatkowych.}
\turn{A}{Kann ich die Wohnung am Donnerstag besichtigen?}{Czy mogę obejrzeć mieszkanie w czwartek?}
\turn{B}{Ja, um achtzehn Uhr. Der Balkon ist im zweiten Stock.}{Tak, o osiemnastej. Balkon jest na drugim piętrze.}
\end{dialog}

\section{Ćwiczenia}
\begin{cwiczenia}
\exitem{Uzupełnij: \de{Der Laptop liegt auf \dots{} Tisch.} (der Tisch, Wo?)}{\de{dem}}
\exitem{Uzupełnij: \de{Ich lege den Laptop auf \dots{} Tisch.} (der Tisch, Wohin?)}{\de{den}}
\exitem{Uzupełnij: \de{Wir sind in \dots{} Küche.} (die Küche, Wo?)}{\de{der}}
\exitem{Uzupełnij: \de{Ich gehe in \dots{} Küche.} (Wohin?)}{\de{die}}
\exitem{Uzupełnij: \de{Das Buch liegt in \dots{} Regal.} (das Regal, Wo?)}{\de{dem}, czyli \de{im}}
\exitem{Uzupełnij: \de{Ich gehe \dots{} Büro.} (Wohin, w skróconej formie)}{\de{ins}}
\exitem{Uzupełnij: \de{Ich bin \dots{} Büro.} (Wo, w skróconej formie)}{\de{im}}
\exitem{Popraw: \de{Ich gehe zu Hause.}}{\de{Ich gehe nach Hause.}}
\exitem{Popraw: \de{Ich bin in das Büro.}}{\de{Ich bin im Büro.}}
\exitem{Wybierz \de{stehen} czy \de{stellen}: \de{Die Flasche \dots{} auf dem Tisch.}}{\de{steht}}
\exitem{Wybierz \de{legen} czy \de{liegen}: \de{Ich \dots{} das Handy auf den Tisch.}}{\de{lege}}
\exitem{Zapytaj o czynsz (formalnie).}{\de{Wie hoch ist die Miete?}}
\exitem{Przetłumacz: „Mieszkanie ma trzy pokoje, kuchnię i balkon.”}{\de{Die Wohnung hat drei Zimmer, eine Küche und einen Balkon.}}
\end{cwiczenia}

\section{Trening słownictwa (generowany)}
\drill{10}

\weekend{Opisz swoje mieszkanie lub wymarzone mieszkanie w Zurychu: pokoje, co gdzie stoi, czynsz.
Użyj przynajmniej pięciu przyimków lokalnych z poprawnym przypadkiem. Poproś AI o sprawdzenie wyłącznie \de{Wo?} / \de{Wohin?}.}
